\documentclass[11pt,letterpaper,twoside]{article}
\usepackage[english]{babel}
\usepackage{amssymb,amsmath}
\usepackage{fancyhdr}
\usepackage{graphicx}
\usepackage{wrapfig}

 \oddsidemargin  0in \evensidemargin 0in
 \topmargin   -0.25in \headheight 0.25in \headsep 0.25in
 \textwidth   6.5in \textheight 8.75in \marginparsep 0pt \marginparwidth 0pt
 \parskip 1ex  \parindent 0ex \footskip 20pt

\newfont{\bssten}{cmssbx10}
\newfont{\bssnine}{cmssbx10 scaled 900}

%%%%%%%%%%%%%%%%%%%%%%%
\newcommand{\whatizit}{Lecture 16: Confidence Intervals for Proportions}
%%%%%%%%%%%%%%%%%%%%%%%


\pagestyle{fancy}  
\fancyhead{\bssnine STOR 155,  \whatizit}
\fancyhead[RE]{} \fancyhead[LO]{}
\fancyhead[LE]{\bssnine \thepage} \fancyhead[RO]{\bssnine \thepage}
\lfoot{} \cfoot{} \rfoot{}   


\newcommand{\var}{\mathrm{Var}}



\begin{document}



\thispagestyle{empty} \vspace*{-0.75in}

{\bssten STOR 155: Introduction to Data Models and Inference \hfill June 12, 2024 \\
Prof. Will Lassiter  \hfill Page 1 of \pageref{totalpag}}
\vspace{10pt}
\begin{center} {{\Large \bf \whatizit}} \end{center}

Motivating example: A drug company is testing ``LDL-Be-Gone", a newly-developed cholesterol medication. The obvious question:  Does it work? \vspace{10pt}

STEP 1: Conduct a well-controlled experiment on 100 volunteers, noting that only 10 of the 50 not using the drug saw a significant decrease in their bad cholesterol levels, while 32 of the 50 using the drug saw a significant decrease. \vspace{10pt}

STEP 2: Spend \$1B mass-producing and marketing the drug? \vspace{120pt}


Recall: Suppose we randomly sample $n$ members of a large population and classify proportion $\hat{p}$ of them as ``successes".  If $p$ represents the proportion of ``successes" in the population at large, the sampling distribution of $\hat{p}$ is ... \vspace{60pt}

Then by the 68-95-99.7 rule, approximately 95\% of sample proportions $\hat{p}$ will fall between ... 

\newpage

Getting more precision with a z-table or Excel: \vspace{200pt}

Definitions: {\bf critical value}, {\bf margin of error}, {\bf confidence interval} \vspace{350pt}

Approximating standard error:

\newpage

Example: We are interested in the proportion $p$ of students at UNC who have engaged in binge drinking behavior at least once within the last month. In a SRS of 485 UNC students, we find that 185 have done so.  

\begin{itemize}

\item Find an (approximate) 95\% confidence interval for $p$. \vspace{120pt}

\item Now suppose your sample size is 4*485=1940 and you have the exact same proportion of binge drinkers.  Construct another 95\% confidence interval based on this larger sample. \vspace{80pt}

\end{itemize}

Interpretation of ``95\% confidence" \vspace{120pt}

Other confidence levels

\newpage

Example: Find $z^*$ for the following confidence levels:

\begin{itemize}

\item 90\% \vspace{50pt}

\item 99\% \vspace{50pt}

\end{itemize}

Example:  The probability of rolling a seven with a pair of fair dice is $\frac{1}{6}$.  We have a pair of dice that may or may not be fair; let $p$ be the unknown probability they will produce a seven.

Suppose we roll this pair of dice 120 times and observe a seven 25 times. 

\begin{itemize}

\item Construct a 95\% confidence interval for $p$ based on these results. \vspace{80pt}

\item Construct a 99\% confidence interval for $p$ based on these results. \vspace{50pt}

\end{itemize}

Shrinking margin of error: 

\newpage

{\bf Choosing a sample size to get desired margin of error} \vspace{6pt}

Example: We plan to take a poll to estimate $p$, the proportion of American voters who support throwing out the current health care system and putting Canada (``America's hat") in charge of US health care.  We want our 95\% margin of error to be $\leq 0.03$.  

\begin{itemize}

\item If we think around 25\% of Americans feel this way, how big of a sample do we need?  \vspace{100pt}

\item What if we have no idea about the value of $p$?

\end{itemize}






\label{totalpag}
\end{document}
